Пусть K - евклидово кольцо \\ \\
\Lemma  У каждого элемента $x \in K\backslash \{0\}$ существует разложение
\\
\Proof Предположим, что это не так и рассмотрим элемент $x \in K\backslash \{0\}$ с минимальной нормой, который не раскладывается в произведение неразложимых. Представим $x$ в виде $x = ab$, где $a$, $b$ $\in K\backslash \{0\}$, тогда $N(x) > N(a)$, $N(x) > N(b)$.
\\
Если оба неравенства строгие, то разложение элемента $x$ существует как произведение разложений элементов $a, b$ -- противоречие. 
\\
Пусть теперь $N(ab) = N(a)$. Разделим $a$ на $ab$ с остатком: $a = q(ab)+r, \;q, r \in K$. 
\begin{itemize}
    \item Если $r = 0$, то $bq = 1$ и $b$ обратим.
    \item  Если же $N(r) < N(ab)$,
то $N(ab) > N(r) = N(a(1 - q b)) > N(a)$,
что невозможно.
\end{itemize}
 Случай, когда $N(ab) = N(b)$, аналогичен.
\\
Таким образом, для всех представлений $x$ в виде $x = ab$, $a$, $b$ $\in K\backslash \{0\}$, один из элементов
$a$, $b$ -- обратимый, поэтому $x$ -- неразложимый элемент, и у него есть разложение $x = 1 \cdot x$.\\
Снова получено противоречие. \EndProof
\bigbreak 
\noindent \Lemma (Евклида) Каждый неразложимый элемент $p$ в $K$ прост, то есть $\forall a, b \in K : p | ab \Longrightarrow p | a $ лии $p | b$.
\\
\Proof Если $p | a$, то утверждение уже доказано. Предположим, что $p \nmid a$,
и докажем, что $\exists x, y \in K : a x + p y = 1$. Для этого рассмотрим следующий процесс последовательного деления с остатком, называемый алгоритмом Евклида:
\begin{itemize}
    \item []$a = p q_1 + r_1$
    \item [] $p = r_1q_2 + r_2$
    \item [] $r1 = r_2 q_3 + r_3$
    \item[] ...
\end{itemize}
Поскольку N($p$) > N($r_1$) > N($r_2$) > . . . , то $\exists k \in \N : r_{k+1} = 0$. Значит, $r_{k-1} = r_k q_{k+1}$, и легко по индукции убедиться, что тогда $r_k | a$ и $r_k | p$. Поскольку p неразложим и $p \nmid a$, то
$r_k \sim 1$. По индукции получим, что $\exists \ x, y \in K : a x + p y = 1$, тогда $abx + pby = b$ и $p | b$.\EndProof
\\
\\
\Th  Кольцо K факториально \\
\Proof В K существует разложение для каждого ненулевого элемента, и каждый неразложимый элемент в K прост по лемме Евклида, поэтому K факториально. \EndProof